\section{Non-Dual Algebra of NDH}
\label{sec:ndh-nondual-algebra}

NDH contains an algebraic structure that governs merge behavior, recursion
symmetry, and fixed-point stability. The algebra is non-dual, meaning that its
operations preserve modal symmetry and do not impose hierarchical or directional
bias. This algebra defines the internal cognitive behavior of NDH.

\subsection{Algebraic Structure}

Let $(\Sigma, \oplus, \odot)$ be an algebra over the cognitive signal space
$\Sigma$. The operations $\oplus$ and $\odot$ satisfy the following axioms:

\begin{enumerate}
\item Commutativity of merge:
\[
x \oplus y = y \oplus x.
\]

\item Distributivity of recursion over merge:
\[
x \odot (y \oplus z) = (x \odot y) \oplus (x \odot z).
\]

\item Idempotence of recursion:
\[
x \odot x = x.
\]
\end{enumerate}

These axioms define NDH's non-dual behavior.

\subsection{Merge Operator}

The merge operator $\oplus : \Sigma \times \Sigma \to \Sigma$ combines spectral
components and recursion outputs. Merge is symmetric and preserves modal
structure. For spectral emission
\[
\mathcal{N}_{DH}(x) = \sum_k \phi_k(x) e_k,
\]
merge acts componentwise:
\[
(\mathcal{N}_{DH}(x) \oplus \mathcal{N}_{DH}(y))
  = \sum_k (\phi_k(x) + \phi_k(y)) e_k.
\]

\subsection{Recursion Operator}

The recursion operator $\odot : \Sigma \times \Sigma \to \Sigma$ defines NDH's
internal recursion behavior. Recursion is modal-preserving and compatible with
holonomy transport. For modal recursion kernels $R_i$,
\[
\mathcal{N}_{DH}(x) \odot \mathcal{N}_{DH}(y)
  = \mathcal{N}_{DH}(R_i(x)).
\]

\subsection{Fixed-Point Behavior}

NDH contains fixed points under recursion:
\[
x \odot x = x.
\]
Fixed points correspond to stable cognitive states within modal regions. For
each modal region $M_i$, the set of fixed points is:
\[
\text{Fix}(R_i) = \{ x \in M_i \mid R_i(x) = x \}.
\]

\subsection{Non-Dual Symmetry}

Non-dual symmetry is defined by the equivalence of modal recursion under merge:
\[
(R_A \oplus R_B) \odot R_C
  = R_C \odot (R_A \oplus R_B).
\]
This symmetry ensures that NDH's internal algebra does not impose directional or
hierarchical structure across modes.

\subsection{Algebra Summary}

The non-dual algebra of NDH consists of a symmetric merge operator, a
modal-preserving recursion operator, fixed-point behavior, and non-dual symmetry
conditions. This algebra defines NDH's internal cognitive behavior and supports
its geometric and spectral structure.
